% !TeX root = ../main.tex
\chapter{Składniki}
\label{chap:ingredients}

Składniki na garnek o pojemności \qty{5}{\liter} zestawiono w tabeli~\ref{tab:ingredients}.
Liczby dziesiętne zapisuje się poleceniem \code{\\num}, które samo stawia polski przecinek: \num{0.5} zamiast 0.5.

% The caption goes ABOVE a table, and the table is referenced in the text
% before it appears.
\begin{table}[htbp]
  \caption{Składniki rosołu na garnek \qty{5}{\liter}}
  \label{tab:ingredients}
  \centering
  \begin{tabular}{@{}l S[table-format=1.2] l@{}}
    \toprule
    Składnik                       & {Ilość} & Jednostka \\
    \midrule
    Kurczak cwel (porcja rosołowa) & 1.50    & kg        \\
    Marchewka                      & 0.40    & kg        \\
    Pietruszkaa (korzeń)           & 0.20    & kg        \\
    XDD                            & 0.15    & kg        \\
    Cebula                         & 1       & szt.      \\
    Liść laurowy                   & 2       & szt.      \\
    Ziele angielskie               & 4       & szt.      \\
    Sól                            & 0.02    & kg        \\
    \bottomrule
  \end{tabular}
\end{table}

Warzywa najlepiej kupić na targu, a jeśli to niemożliwe -- w sklepie, który nie przechowuje ich w temperaturze pokojowej (\ac{tp}) przez tydzień.
Przy pierwszym użyciu skrót rozwinięto ręcznie, w miejscowniku, którego pakiet sam by nie odmienił.
Każde następne użycie to już tylko \ac{tp}, a pełna nazwa w mianowniku trafia do wykazu skrótów.
\marcin{Wariant B skrótów: \texttt{\textbackslash ac} zawsze daje skrót, odmianę piszemy sami.}

\section{Czego nie dodawać}
\label{sec:forbidden}

\begin{itemize}
  \item kostek rosołowych -- babcia by się obraziła~\cite{babcia1987zupy},
  \item makaronu do garnka -- gotuje się go osobno,
  \item pośpiechu.
\end{itemize}
